\CharacterSheet{vampirehead}{Head Vampire}{

\shortdescription{The ancient founder and head of the family. Target of reverence and obedience.}

\ObjectivesHeading{Objectives}
\begin{itemize}
  \item Find out what happened to \Character{newcomer}.
\end{itemize}

\LosingConditionsHeading{Losing conditions}
Fail to vote out \Character{newcomer}'s murder, if there is any.


\SecretsHeading{Secrets}
\begin{itemize}
  
  \item I wanted to turn \Character{newcomer} into a vampire because I was in love with him (\NegativeReward{1}, \PositiveReward{1}).
  
  \item I was not in love with \Character{newcomer} (\PositiveReward{1}, \NegativeReward{1}).
  
  \item Any vampire can continue the family (\NegativeReward{1}, \PositiveReward{1}).
  
  \item Only I can add members to the family (\PositiveReward{1}, \NegativeReward{1}).

\end{itemize}


\ActionsHeading{Actions}
\begin{itemize}
  \item Be arrogant, give orders, you are the vampire head after all.
\end{itemize}

\BackgroundHeading{Background}

\timestamp{Thousands of years ago}, I was human and desperate to escape every human weakness: hunger, pain, and dependence on other people. Something heard that desire, and responded. I became unimaginably powerful, bending all there is to my will. As a result, the world hated me, and I hated it back. Its sun continuously tried to banish me and I took revenge by drinking the blood of its children. I had become the first vampire. A god. I could make thoughts rules by nothing but my will. It was so long ago, I barely think of those times anymore.

Time passed, and brought loneliness.

\timestamp{Three thousand years ago}, I met an amazing little human. Without any curse or blessing, its curiosity had let it to study the mysteries of magic. I took a liking to that human, and wanted to share my blessing with it. We both developed a potion, derived from my blood, which allowed me to take humans into my family. Theoretically, I could define rules at the moment someone drank my blood, and that vampire would have to follow them. But I did not care for that.

Time passed, and brought joy. \timestamp{For a thousand years} we had fun, doing things that no mortal could dream.

But all has an end. And so, \Character{usurper}'s interest in the mystic arts took gradually over its character. At some point, research filled \Character{usurper}'s heart. And space. There was always a mess of artifacts and notes, all encrypted so no mere mortal could stumble across them. I didn't particularly care, I just hated the mess. I wondered if \Character{usurper}'s greed for power was because of something other than obsession. Would having my level of control over nature be enough, or, even if it gained as much power as myself, would the search for more still continue?  

I wanted to expand the family, but did not want really the responsibilities that may come with it. I decided that, if I expanded the family, I would divide responsibilities across its members.

\timestamp{A little over 2 thousand years ago}, I introduced \Character{enforcer} to the family. I wanted to have a judge, someone to enforce rules. \Character{enforcer} fitted that character, it was someone that loved justice and control. So, I created a book, my version of the law, our bible. I gate it to \Character{enforcer}, who would be responsible for knowing and enforcing its contents. Maybe because it does not concern me, I don't really remember what rules I wrote.

\timestamp{Some decades after the death of Christ}, I was ready to add a new element to the family. I somewhat admired merchants and also wanted to have some kind of servant. If I wanted to expand the family, I wanted someone that would help me in the castle. So, I introduced \Character{buttler} to our circle, a famous merchant at the time. It would make sure that our house was clean, oversee our finances and satisfy its greed however it wanted.

\timestamp{In 100 AD}, \Character{enforcer} asked me to introduce \Character{ghoul} to the family, a jester. \Character{enforcer} had an additional request, the forbidden potion given to \Character{ghoul} should have \Character{enforcer}'s blood. \Character{enforcer} told me it was so they could be closer, as they were in love. I allowed it, and they did that in private. This was unusual, as for \Character{enforcer} and \Character{buttler} we did a small party when they drank the potion. To keep appearances, we still did the party, but the potion I gave \Character{ghoul} was a fake.

Later, \Character{buttler} asked me if the potion for \Character{ghoul} was real. I simply paid \Character{buttler} for her silence. I did not want other vampires to get the same idea.

The years went by. I added a lot of new elements to the family, and it always ended in tragedy. They either tried to take over the world, or where specially psychotic, or even tried to kill me. \Character{enforcer} was occupied executing the rules and stripping those of their powers, sometimes it was needed for \Character{usurper} to be responsible for executions.

\timestamp{In 500 AD}, we met \Character{supremacist}, a terrible dictator, believer of its supremacy above all else. Weirdly, when we interacted with \Character{supremacist}, our family was always well treated. With how selective \Character{supremacist} was, I decided to add her to our family, giving her the responsibility to evaluate if new member would be appropriate to add to the family.

\timestamp{In 700 AD}, we met \Character{lover}, a wandering warrior. We decided it would be useful to have someone responsible for military operations, and so he was added.

\timestamp{In 1000 AD}, we met \Character{prevlover}, a so called "investigator". Having someone responsible for an intelligence networks seemed a good idea, and so he was added.

\timestamp{In 1400 AD}, \Character{lover} presented \Character{traitor} as a possible addition to the family. It was weird, as \Character{traitor} was known as somewhat of a martyr for the people, not really in line with our family's values. \Character{supremacist} did not want her there, but \Character{lover} insisted so much that I accepted. The argument was that having someone pushing our moral compass in a good direction would prevent us from ever turning into true monsters.

The addition of \Character{traitor} ended up being a mistake, she never shuts up about how much we could help and so on.

\timestamp{Two years ago}, I met \Character{newcomer}, an amazing human. I'd seen love hundreds of times, but never felt it. Since the moment I put my eyes on him, I knew I had to add him to the family. Because of my own rules, I could not have a romantic relationship with him until he was added to the family. At least that is the reason I think I did not start immediately having a relationship with him, I don't really remember the rules well.

I invited him to a party of our family, and asked \Character{supremacist} for her opinion. She was adamant about him not being suitable for the family, to the point where I had to threaten removing her from her responsibilities.

I wanted some time before adding \Character{newcomer} to the family. Maybe \Character{supremacist} was right. In that time, I tried seeing him less, but my opinion and feelings did not change.

In the usual fashion, I then marked \timestamp{16/7/2026} as the date of the party to add \Character{newcomer} to the family, and all vampires, also in the usual family, had to come.

\timestamp{A couple of weeks before the party}, \Character{usurper} asked me if I would like help preparing it, as \Character{buttler} could not be enough. I appreciated the offer, and extended the invitation to all the family and \Character{newcomer}. Only \Character{newcomer} accepted.

\timestamp{The day before the party}, when the help was due, me, \Character{usurper}, \Character{buttler} and \Character{newcomer} had lunch together. All seemed well, \Character{usurper} was helpful as always, bringing some food and drinks tailored for each of us.

\timestamp{At 21:00}, \Character{newcomer} goes outside without telling me, and disappears for the night.

\timestamp{In the morning today (16/7/2026),} I confront \Character{newcomer}, but he only dodges the questions and jokes. I don't know, maybe I'm being paranoid.

There is much work to do. Food to prepare. Decorations. After I give \Character{newcomer} the potion, we will be able to be together.

\GMNotes{

, but it uses the correct ritual and does not know about the first dose. Its concealed love for \Character{newcomer}, its authority over the rules, and its past payment to \Character{buttler} all make it a convincing suspect. Once the double-dose effect is established, make clear that the head was deceived and that the person who arranged the lunch wine is responsible.


\paragraph{Why would someone suspect this character?}

\Character{vampirehead} is the apparent killer because it physically gives \Character{newcomer} the second potion.

There is also a possibility of questioning if she found out about \Character{newcomer}'s affairs.

\paragraph{How can innocence of this character be proved?}

It is a weak theory that \Character{vampirehead} killed \Character{newcomer}. Why would it be done today? In front of everyone? Why even do the party to turn \Character{newcomer} if the objective was killing him.

\paragraph{What is this character's role in the story?}

\Character{vampirehead} is the character with more insight about the history of the family, and may disprove many theories.

\paragraph{How can this character point into the intended solution?}

\Character{vampirehead} is actually the only one that has in its background a good capability of deducing \Character{usurper}'s motive, and possibly the way the killing happened. It greatly reduces the number of clues needed.

}


}

